% Generated draft: requires mathematical and source review.
\documentclass[11pt]{article}
\usepackage[margin=25mm]{geometry}
\usepackage{fontspec}
\setmainfont{Latin Modern Roman}
\usepackage{amsmath,amssymb,mathtools,unicode-math}
\setmathfont{Latin Modern Math}
\usepackage{enumitem,needspace,xurl,hyperref}
\setlength{\parindent}{0pt}
\setlength{\parskip}{5pt}
\setlength{\emergencystretch}{3em}
\newcommand{\N}{\mathbb{N}}
\newcommand{\Z}{\mathbb{Z}}
\newcommand{\Q}{\mathbb{Q}}
\newcommand{\R}{\mathbb{R}}
\newcommand{\C}{\mathbb{C}}
\newcommand{\F}{\mathbb{F}}
\newcommand{\A}{\mathbb{A}}
\newcommand{\PP}{\mathbb{P}}
\newcommand{\E}{\mathbb{E}}
\newcommand{\eps}{\varepsilon}
\newcommand{\dd}{\,\mathrm{d}}
\begin{document}
\section*{UnsolvedMath ID: 10600045}
\textbf{Source problem code:} catalogue entry
\subsection{Definitions and mathematical statement}
For an integer \(n\ge3\), let \(VB_n\) be the virtual braid group on \(n\) strands: diagrams have classical and virtual crossings, are monotone between prescribed top and bottom endpoints, and are identified by the virtual braid moves fixing that boundary data. Let \(\mathcal F_n\) denote V. O. Manturov's invariant \(\mathcal F\) for such braids. Its construction, target and equality relation on values are not supplied in the record.

Problem-list item 45 asks to prove completeness for every strand number greater than two:
\[
 \forall n\ge3\ \forall\beta,\gamma\in VB_n:\qquad
 \mathcal F_n(\beta)=\mathcal F_n(\gamma)
 \ \Longrightarrow\ \beta=\gamma\text{ in }VB_n.
\]
The reverse implication is the stated invariance property. This is a question about distinguishing braids with their fixed endpoints, not about equivalence of braid closures under Markov moves. The source specifies no exception at a particular \(n\ge3\), but it also gives no claim for one or two strands. Without the definition of \(\mathcal F\), this extract cannot establish a precise target criterion beyond the injectivity statement. In particular it does not justify replacing \(\mathcal F\) by a different known virtual-braid representation.
\subsection{Short English statement}
Prove that Manturov's specified invariant distinguishes every pair of inequivalent virtual braids on each number of strands greater than two. Its defining construction is missing, so completeness remains an exact injectivity request attached to an unresolved source reference.
\subsection{Sources}
\begin{itemize}
\item Local UnsolvedMath registry, permanent ID \(10600045\), exact source code \texttt{catalogue entry}, virtual-knot problem-list item 45. The invariant is attributed to V. O. Manturov; no proposed year is supplied. The title is truncated.
\item \emph{Unsolved Problems in Virtual Knot Theory and Combinatorial Knot Theory}, revised survey, dedicated problem-list item 45: \url{https://arxiv.org/abs/1409.2823}. Background: 60 top-level items, compound items retained intact; source-TeX extraction, accessed 2026-07-29.
\item Supplied related reference: ``Representations of virtual braids by automorphisms and virtual knot groups'' (2016), arXiv:1603.xxxx. This is an incomplete placeholder identifier and is not a supplied definition of \(\mathcal F\).
\item Source evidence NEEDS\_REVIEW presents an open problem and requires a current primary-literature audit. Supplied OPEN-TRIAGE summary reports no resolution found and describes completeness as apparently open; no independent verification is asserted.
\item Dataset attribution: Ulam AI / UnsolvedMath Contributors, CC BY 4.0. Source rights NEEDS\_REVIEW: publicly accessible author/publisher copy, redistribution terms requiring review.
\end{itemize}
\end{document}
